\section{Related Work}
\label{sec:related_work}

\textbf{Parameter-Efficient Fine-Tuning and Model Merging.} 
The proliferation of large, pre-trained base models has accelerated the development of Parameter-Efficient Fine-Tuning (PEFT) methods, particularly Low-Rank Adaptation (LoRA) \cite{hu2021lora}, residual adapters \cite{houlsby2019parameter, rebuffi2017learning}, prompt tuning \cite{lester2021power}, and prefix tuning \cite{li2021prefix}. Fusing these specialized models at runtime has emerged as a cost-effective alternative to maintaining multiple separate instances. Research in model merging typically focuses on static parameter-space combinations, such as model soups \cite{wortsman2022soups}, task arithmetic \cite{ilharco2022editing}, or projection into tangent space \cite{ortiz2023task}. While computationally efficient, static weight-space merging ignores inputs at runtime and is prone to interference (representation cross-talk) when merging highly disparate tasks. 

\textbf{Dynamic Gating and Mixture of Experts.} 
Sparsely Gated Mixture of Experts (MoE) \cite{shazeer2017outrageously} and routing-based architectures \cite{fedus2022switch} resolve task interference by evaluating routing gates dynamically at runtime. Traditional MoE approaches utilize routing heads to dynamically partition the input space and activate task-specific experts \cite{jacobs1991adaptive, jordan1994hierarchical}. In sequential deep networks, layer-by-layer dynamic activation ensembling has been proposed to optimize multi-task serving environments \cite{sable2025samplewise}. However, these sequential routers are notorious for high transductive overfitting under extreme data scarcity and often exhibit violent gating variations across network depth. 

\textbf{Trajectory Smoothing and Model Kinetics.} 
To stabilize routing paths across sequential layers, recent works have proposed trajectory-level heuristics. For example, \cite{chemmerge2025kinetics} introduces continuous kinetics differential equations to temporalize routing coefficients, mimicking reaction rate dynamics. Similarly, SABLE \cite{sable2025samplewise} utilizes non-parametric nearest-centroid projections, and RBPM \cite{rbpm2026rademacher} constrains gating coefficients to low-degree polynomials across layers. Although these methods improve empirical smoothness, they operate as heuristic approximations and lack formal convergence or stability proofs. Our work fills this critical theoretical void.

\textbf{Lipschitz Constraints and Spectral Normalization.} 
Regularizing the Lipschitz constant of deep neural networks is a well-established technique to guarantee generalization, robustness, and stability. Pioneering works have leveraged spectral normalization on weight matrices to control representation drift and ensure smooth function mapping in generative networks \cite{miyato2018spectral}, defense against adversarial perturbations \cite{tsuzuku2018lipschitz}, and regularizing general architectures \cite{gouk2021regularisation, anil2019sorting, bartlett2017spectrally}. Formulations analyzing deep networks as dynamical systems often study continuous limits or contraction mappings to guarantee stability \cite{hale1969ordinary, scaman2018lipschitz, virmaux2018lipschitz}. Drawing on functional analysis \cite{rudin1976principles, kreyszig1978introductory} and Banach's Fixed-Point Theorem \cite{banach1922operations}, we establish the first explicit contraction mapping framework for deep sequential model ensembling.
